% Bagian: turing-machines
% Bab: undecidability
% Subbagian: verification

\documentclass[../../../include/open-logic-section]{subfiles}

\begin{document}

\olfileid[id]{tur}{und}{ver}
\olsection{Memverifikasi Representasi}

\begin{explain}
Untuk memverifikasi bahwa representasi kita bekerja, kita harus membuktikan
dua hal. Pertama, kita harus menunjukkan bahwa jika $M$ berhenti pada
masukan~$w$, maka $!T(M, w) \lif !E(M, w)$ valid. Kemudian, kita harus
menunjukkan arah sebaliknya, yakni bahwa jika $!T(M, w) \lif !E(M, w)$ valid,
maka $M$ memang akhirnya berhenti ketika dijalankan pada masukan~$w$.

Strategi untuk membuktikan kedua hal ini sangat berbeda. Untuk hasil pertama,
kita harus menunjukkan bahwa !!a{sentence} logika orde pertama (yaitu
$!T(M, w) \lif !E(M, w)$) valid. Cara termudah untuk melakukannya ialah
memberikan !!a{derivation}. Namun, bukti kita harus berlaku bagi semua $M$
dan~$w$, sehingga sebenarnya bukan hanya ada satu !!{sentence} yang perlu kita
beri !!a{derivation}, melainkan tak berhingga banyaknya. Karena itu, yang dapat
kita lakukan ialah membuktikan dengan induksi bahwa bagaimana pun bentuk $M$
dan~$w$, dan berapa pun banyak langkah yang diperlukan $M$ untuk berhenti pada
masukan~$w$, akan ada !!a{derivation} bagi $!T(M, w) \lif !E(M, w)$.

Secara alami, induksi kita berjalan pada banyaknya langkah yang diambil $M$
sebelum mencapai konfigurasi penghentian. Dalam bukti induktif, kita akan
menetapkan bahwa untuk setiap langkah~$n$ dalam jalannya $M$ pada masukan~$w$,
$!T(M, w) \Entails !C(M, w, n)$, dengan $!C(M, w, n)$ mendeskripsikan secara
tepat konfigurasi $M$ yang dijalankan pada~$w$ setelah~$n$ langkah. Jika $M$
berhenti pada masukan~$w$ setelah, katakanlah, $n$ langkah, maka $!C(M, w, n)$
mendeskripsikan konfigurasi penghentian. Kita juga akan menunjukkan bahwa
$!C(M, w, n) \Entails !E(M, w)$ setiap kali $!C(M, w, n)$ mendeskripsikan
konfigurasi penghentian. Jadi, jika $M$ berhenti pada masukan~$w$, maka untuk
suatu~$n$, $M$ berada dalam konfigurasi penghentian setelah $n$ langkah.
Dengan demikian, $!T(M, w) \Entails !C(M, w, n)$, dengan $!C(M, w, n)$
mendeskripsikan konfigurasi penghentian; dan karena dalam kasus itu
$!C(M, w, n) \Entails !E(M, w)$, kita memperoleh
$!T(M, w) \Entails !E(M, w)$, yakni $\Entails !T(M, w) \lif !E(M, w)$.

Strategi untuk arah sebaliknya sangat berbeda. Di sini kita mengandaikan
$\Entails !T(M, w) \lif !E(M, w)$ dan harus membuktikan bahwa $M$ berhenti
pada masukan~$w$. Dari hipotesis diperoleh $!T(M, w) \Entails !E(M, w)$,
yakni $!E(M, w)$ benar dalam setiap !!{structure} tempat $!T(M, w)$ benar.
Karena itu, kita akan mendeskripsikan !!a{structure}~$\Struct{M}$ tempat
$!T(M, w)$ benar: domainnya adalah $\Nat$, sedangkan interpretasi semua
$\Obj Q_q$ dan $\Obj S_\sigma$ diberikan oleh konfigurasi-konfigurasi~$M$
selama dijalankan pada masukan~$w$. Sebagai contoh,
$\Sat{M}{\Obj Q_q(\num{m}, \num{n})}$ berlaku jika dan hanya jika $M$, ketika
dijalankan pada masukan~$w$ selama $n$ langkah, berada dalam keadaan~$q$ dan
memindai petak~$m$. Karena berdasarkan hipotesis
$!T(M, w) \Entails !E(M, w)$, dan berdasarkan konstruksi
$\Sat{M}{!T(M, w)}$, maka $\Sat{M}{!E(M, w)}$. Namun,
$\Sat{M}{!E(M, w)}$ berlaku jika dan hanya jika terdapat suatu
$n \in \Domain{M} = \Nat$ sehingga $M$, ketika dijalankan pada masukan~$w$,
berada dalam konfigurasi penghentian setelah~$n$ langkah.
\end{explain}

\begin{defn}
Misalkan $!C(M, w, n)$ adalah !!{sentence}
\[
\Obj Q_q(\num{m}, \num{n}) \land \Obj S_{\sigma_0}(\num{0}, \num{n})
\land \dots \land \Obj S_{\sigma_k}(\num{k}, \num{n}) \land
\lforall[x][(\num{k} < x \lif \Obj S_\TMblank(x, \num{n}))]
\]
di mana $q$ adalah keadaan $M$ pada saat~$n$, $M$ memindai petak~$m$ pada
saat~$n$, petak~$i$ memuat simbol~$\sigma_i$ pada saat~$n$ untuk
$0 \le i \le k$, dan $k$ adalah yang lebih besar di antara indeks petak
tak-kosong paling kanan pada pita ketika saat~$0$ dan indeks petak paling kanan
yang pernah dikunjungi kepala pita setelah $n$ langkah.
\end{defn}

\begin{lem}\ollabel{lem:halt-config-implies-halt}
Jika $M$ yang dijalankan pada masukan~$w$ berada dalam konfigurasi penghentian
setelah $n$ langkah, maka $!C(M, w, n) \Entails !E(M, w)$.
\end{lem}

\begin{proof}
Andaikan $M$ berhenti pada masukan~$w$ setelah $n$ langkah. Terdapat suatu
keadaan~$q$, petak~$m$, dan simbol~$\sigma$ sedemikian sehingga:
\begin{enumerate}
\item Setelah $n$ langkah, $M$ berada dalam keadaan~$q$ dan memindai
  petak~$m$ yang memuat~$\sigma$.
\item Fungsi transisi $\delta(q, \sigma)$ tidak terdefinisi.
\end{enumerate}
$!C(M, w, n)$ adalah deskripsi konfigurasi ini dan memuat klausa
$\Obj Q_{q}(\num{m}, \num{n})$ serta
$\Obj S_{\sigma}(\num{m}, \num{n})$. Kedua klausa ini bersama-sama
menyiratkan $!E(M,w)$:
\[
\lexists[x][\lexists[y][(\bigvee_{\tuple{q, \sigma} \in
      X}(\Obj Q_q(x, y) \land \Obj S_\sigma(x, y)))]]
\]
sebab $\Obj Q_q(\num{m}, \num{n}) \land
\Obj S_\sigma(\num{m}, \num{n}) \Entails
\bigvee_{\tuple{q', \sigma'} \in X}
(\Obj Q_{q'}(\num{m}, \num{n}) \land
\Obj S_{\sigma'}(\num{m}, \num{n}))$, karena
$\tuple{q,\sigma} \in X$.
\end{proof}

\begin{explain}
Jadi, jika $M$ berhenti pada masukan~$w$, terdapat suatu~$n$ sedemikian sehingga
$!C(M, w, n) \Entails !E(M,w)$. Sekarang akan kita tunjukkan bahwa untuk
setiap saat~$n$, $!T(M, w) \Entails !C(M, w, n)$.
\end{explain}

\begin{lem}
\ollabel{lem:config}
Untuk setiap $n$ sedemikian sehingga jalannya $M$ pada masukan~$w$ memiliki
konfigurasi setelah $n$ langkah, $!T(M, w) \Entails !C(M, w, n)$.
\end{lem}

\begin{proof}
Basis induksi: Jika $n = 0$, maka konjungsi-konjungsi $!C(M, w, 0)$ juga
merupakan konjungsi-konjungsi $!T(M, w)$, sehingga diperikutkan olehnya.

Hipotesis induksi: jika jalannya mesin telah mencapai konfigurasi setelah
langkah ke-$n$, maka $!T(M,w) \Entails !C(M, w, n)$. Kita harus menunjukkan
bahwa, kecuali $!C(M, w, n)$ mendeskripsikan konfigurasi penghentian,
$!T(M, w) \Entails !C(M, w, n+1)$.

Andaikan setelah $n$ langkah, $M$ yang dijalankan pada~$w$ berada dalam
keadaan~$q$ dan memindai petak~$m$. Karena $M$ tidak berhenti setelah~$n$
langkah, program~$M$ pasti memuat instruksi dalam salah satu dari tiga bentuk
berikut:

\begin{enumerate}
\item \ollabel{right} $\delta(q, \sigma) = \tuple{q', \sigma', \TMright}$

\item \ollabel{left} $\delta(q, \sigma) = \tuple{q', \sigma', \TMleft}$

\item \ollabel{stay} $\delta(q, \sigma) = \tuple{q', \sigma', \TMstay}$
\end{enumerate}

Ketiga kasus ini akan kita bahas satu per satu.

\begin{enumerate}
\item Andaikan terdapat instruksi dalam bentuk~\olref{right}.
  Menurut \olref[rep]{defn:tm-descr}\olref[rep]{rep-right}, ini berarti
\begin{align*}
& \lforall[x][\lforall[y][((\Obj Q_{q}(x,
  y) \land \Obj S_\sigma(x, y)) \lif {}]]\\
& \qquad (\Obj Q_{q'}(x',
  y') \land \Obj S_{\sigma'}(x, y') \land !A(x, y)))
\intertext{merupakan konjungsi $!T(M,w)$. Ini memperikutkan !!{sentence}
  berikut (instansiasi universal, $\num{m}$ untuk~$x$ dan
  $\num{n}$ untuk~$y$):}
& (\Obj Q_{q}(\num{m}, \num{n}) \land \Obj S_{\sigma}(\num{m},
\num{n})) \lif {}\\
& \qquad (\Obj Q_{q'}(\num{m}', \num{n}') \land
\Obj S_{\sigma'}(\num{m}, \num{n}') \land !A(\num{m}, \num{n})).
\intertext{Menurut hipotesis induksi, $!T(M, w) \Entails !C(M, w, n)$,
  yakni}
& \Obj Q_q(\num{m}, \num{n}) \land \Obj S_{\sigma_0}(\num{0}, \num{n})
\land \dots \land \Obj S_{\sigma_k}(\num{k}, \num{n}) \land \\
& \qquad\lforall[x][(\num{k} < x \lif \Obj S_\TMblank(x, \num{n}))]\\
\intertext{Karena setelah $n$ langkah petak pita~$m$ memuat~$\sigma$,
  konjungsi yang bersesuaian adalah~$\Obj S_\sigma(\num{m}, \num{n})$,
  sehingga ini memperikutkan:}
& \Obj Q_{q}(\num{m}, \num{n}) \land \Obj S_{\sigma}(\num{m},
   \num{n}).
\intertext{Sekarang kita memperoleh}
& \Obj Q_{q'}(\num{m}', \num{n}') \land \Obj S_{\sigma'}(\num{m},
  \num{n}') \land {}\\
& \qquad \bigwedge_{\substack{0\le i\le k\\i\ne m}}
  \Obj S_{\sigma_i}(\num{i}, \num{n}') \land {}\\
& \qquad \lforall[x][(\num{k} < x \lif \Obj S_\TMblank(x, \num{n}'))].
\end{align*}
Baris pertama diperoleh langsung dari konsekuen implikasi sebelumnya melalui
modus ponens. Setiap konjungsi di baris tengah---yang tidak mencakup
$\Obj S_{\sigma_m}(\num{m},\num{n}')$---diperoleh dari konjungsi yang
bersesuaian dalam~$!C(M, w, n)$ bersama dengan
$!A(\num{m},\num{n})$.

Jika $m < k$, maka $!T(M,w) \Proves \num{m} < \num{k}$
(\olref[rep]{prop:mlessk}); berdasarkan transitivitas~$<$, kita memperoleh
$\lforall[x][(\num{k} < x \lif \num{m} < x)]$. Jika $m = k$, maka
$\lforall[x][(\num{k} < x \lif \num{m} < x)]$ berdasarkan logika saja.
Baris terakhir kemudian diperoleh dari konjungsi terakhir dalam
$!C(M,w,n)$, implikasi tersebut, dan $!A(\num{m},\num{n})$.
Jika $m<k$, hasil ini sudah merupakan $!C(M, w, n+1)$.

Sekarang andaikan $m=k$. Dalam kasus ini, setelah $n+1$ langkah kepala pita
juga telah mengunjungi petak~$k+1$, yang kini merupakan petak paling kanan
yang pernah dikunjungi. Karena itu, $!C(M, w, n+1)$ memiliki konjungsi baru
$\Obj S_\TMblank(\num{k}',\num{n}')$, dan konjungsi terakhirnya adalah
$\lforall[x][(\num{k}' < x \lif \Obj S_\TMblank(x, \num{n}'))]$.
Kita harus memverifikasi bahwa kedua !!{sentence} ini juga diperikutkan.

Kita sudah memiliki $\lforall[x][(\num{k} < x \lif \Obj S_\TMblank(x,
\num{n}'))]$. Secara khusus, ini memberi kita
$\num{k} < \num{k}' \lif \Obj S_\TMblank(\num{k}', \num{n}')$.
Dari aksioma $\lforall[x][x < x']$ diperoleh $\num{k} < \num{k}'$.
Melalui modus ponens, diperoleh $\Obj S_\TMblank(\num{k}',\num{n}')$.

Selain itu, karena $!T(M,w) \Proves \num{k} < \num{k}'$, aksioma
transitivitas~$<$ bersama dengan konjungsi pita-kosong di atas memberi kita
$\lforall[x][(\num{k}' < x \lif \Obj S_\TMblank(x, \num{n}'))]$.
(Verifikasinya kita serahkan sebagai latihan.)

\item Andaikan terdapat instruksi dalam bentuk~\olref{left}.
  Menurut \olref[rep]{defn:tm-descr}\olref[rep]{rep-left},
\begin{align*}
& \lforall[x][\lforall[y][((\Obj Q_{q}(x', y) \land \Obj
    S_{\sigma}(x', y)) \lif {}]]\\
& \qquad (\Obj Q_{q'}(x, y') \land \Obj
  S_{\sigma'}(x', y') \land !A(x', y))) \land {}\\
& \lforall[y][((\Obj Q_q(\num{0}, y) \land \Obj S_{\sigma}(\num{0},
    y)) \lif {}]\\
& \qquad (\Obj Q_{q'}(\num{0}, y') \land \Obj S_{\sigma'}(\num{0}, y')
  \land !A(\num{0}, y)))
\intertext{merupakan konjungsi $!T(M,w)$. Jika $m>0$, misalkan $l=m-1$
  (yakni $m=l+1$). Konjungsi pertama !!{sentence} di atas memperikutkan:}
& (\Obj Q_{q}(\num{l}', \num{n}) \land \Obj S_{\sigma}(\num{l}', \num{n}))
\lif {} \\
& \qquad (\Obj Q_{q'}(\num{l}, \num{n}') \land \Obj S_{\sigma'}(\num{l}',
\num{n}') \land !A(\num{l}', \num{n})).
\intertext{Jika tidak, misalkan $l=m=0$ dan perhatikan !!{sentence} berikut
  yang diperikutkan oleh konjungsi kedua:}
& ((\Obj Q_q(\num{0}, \num{n}) \land \Obj S_{\sigma}(\num{0}, \num{n})) \lif {}\\
& \qquad (\Obj Q_{q'}(\num{0}, \num{n}') \land \Obj S_{\sigma'}(\num{0}, \num{n}') \land
!A(\num{0}, \num{n}))).
\intertext{Kedua kalimat itu memperikutkan}
& \Obj Q_{q'}(\num{l}, \num{n}') \land \Obj S_{\sigma'}(\num{m},
  \num{n}') \land {}\\
&\qquad \bigwedge_{\substack{0\le i\le k\\i\ne m}}
  \Obj S_{\sigma_i}(\num{i}, \num{n}') \land {} \\
& \qquad \lforall[x][(\num{k} < x
  \lif \Obj S_\TMblank(x, \num{n}'))].
\end{align*}
seperti sebelumnya. (Perhatikan bahwa dalam kasus pertama,
$\num{l}' \ident \num{l+1} \ident \num{m}$, dan dalam kasus kedua
$\num{l} \ident \num{0}$.) Namun, hasil ini tepat merupakan
$!C(M, w, n+1)$.

\item Kasus \olref{stay} diserahkan sebagai latihan.
\end{enumerate}
Kita telah menunjukkan bahwa bagi setiap~$n$ yang dicapai oleh jalannya mesin,
$!T(M, w) \Entails !C(M, w, n)$.
\end{proof}

\begin{prob}
Lengkapilah kasus~\olref[tur][und][ver]{stay} dalam bukti
\olref[tur][und][ver]{lem:config}.
\end{prob}

\begin{prob}
Berikan !!a{derivation} bagi $\Obj S_{\sigma_i}(\num{i}, \num{n}')$ dari
$\Obj S_{\sigma_i}(\num{i}, \num{n})$ dan
$!A(\num{m}, \num{n})$ (dengan mengandaikan $i \neq m$, yakni $i < m$ atau
$m < i$).
\end{prob}

\begin{prob}
Berikan !!a{derivation} bagi $\lforall[x][(\num{k}' < x \lif \Obj
  S_\TMblank(x, \num{n}'))]$ dari $\lforall[x][(\num{k} < x \lif \Obj
  S_\TMblank(x, \num{n}'))]$, $\lforall[x][x < x']$, dan
$\lforall[x][\lforall[y][\lforall[z][ ((x < y \land y < z) \lif x <
      z)]]]$.
\end{prob}


\begin{lem}
\ollabel{lem:valid-if-halt}
Jika $M$ berhenti pada masukan~$w$, maka $!T(M, w) \lif
!E(M, w)$ valid.
\end{lem}

\begin{proof}
Menurut \olref{lem:config}, bagi setiap saat~$n$ yang dicapai oleh jalannya
mesin, deskripsi~$!C(M, w, n)$ atas konfigurasi $M$ pada saat~$n$
diperikutkan oleh~$!T(M, w)$. Andaikan $M$ berhenti setelah $k$ langkah.
Pada saat itu mesin memindai suatu petak~$m$, dengan $m \in \Nat$.
Maka $!C(M, w, k)$ mendeskripsikan konfigurasi penghentian~$M$, yakni memuat
sebagai konjungsi baik $\Obj Q_q(\num{m}, \num{k})$ maupun
$\Obj S_\sigma(\num{m}, \num{k})$, dengan $\delta(q,\sigma)$ tidak
terdefinisi. Jadi, menurut \olref{lem:halt-config-implies-halt},
$!C(M, w, k) \Entails !E(M,w)$. Namun, karena
$!T(M, w) \Entails !C(M, w, k)$, kita memperoleh
$!T(M, w) \Entails !E(M, w)$; dengan demikian,
$!T(M, w) \lif !E(M, w)$ valid.
\end{proof}


\begin{explain}
Untuk menuntaskan verifikasi klaim kita, kita juga harus menetapkan arah
sebaliknya: jika $!T(M, w) \lif !E(M, w)$ valid, maka $M$ memang berhenti
ketika dijalankan pada masukan~$w$.
\end{explain}

\begin{lem}
\ollabel{lem:halt-if-valid}
Jika $\Entails !T(M, w) \lif !E(M, w)$, maka $M$ berhenti pada masukan~$w$.
\end{lem}

\begin{proof}
Perhatikan $\Lang L_M$-!!{structure}~$\Struct M$ dengan domain~$\Nat$ yang
menafsirkan $\Obj 0$ sebagai~$0$, $\prime$ sebagai fungsi penerus, dan $<$
sebagai relasi lebih-kecil-daripada, serta menafsirkan predikat $\Obj Q_q$
dan~$\Obj S_\sigma$ sebagai berikut:
\begin{align*}
  \Assign{\Obj Q_q}{M} & =
\Setabs{\tuple{m, n}}{\begin{array}{ll}\text{setelah dijalankan pada~$w$ selama~$n$ langkah,}\\ \text{$M$ berada dalam keadaan $q$
  dan memindai petak~$m$}\end{array}} \\
\Assign{\Obj S_\sigma}{M} & = \Setabs{\tuple{m, n}}{\begin{array}{ll}
\text{setelah dijalankan pada~$w$ selama $n$ langkah,}\\ \text{petak~$m$ pada $M$ memuat
  simbol~$\sigma$}\end{array}}.
\end{align*}
Dengan kata lain, kita mengonstruksi !!{structure}~$\Struct{M}$ agar
mendeskripsikan, langkah demi langkah, apa yang sebenarnya dilakukan $M$
ketika dijalankan pada masukan~$w$. Jelas bahwa $\Sat{M}{!T(M, w)}$.
Jika $\Entails !T(M, w) \lif !E(M, w)$, maka $\Sat{M}{!E(M, w)}$ pula, yakni
\[
\Sat{M}{\lexists[x][\lexists[y][(\bigvee_{\tuple{q, \sigma} \in
      X}(\Obj Q_q(x, y) \land \Obj S_\sigma(x, y)))]]}.
\]
Karena $\Domain{M} = \Nat$, pasti terdapat $m,n \in \Nat$ sehingga
$\Sat{M}{\Obj Q_q(\num{m}, \num{n}) \land \Obj S_\sigma(\num{m},
\num{n})}$ bagi suatu~$q$ dan~$\sigma$ dengan $\delta(q, \sigma)$ tidak
terdefinisi. Berdasarkan definisi~$\Struct M$, ini berarti bahwa setelah
dijalankan pada masukan~$w$ selama~$n$ langkah, $M$ berada dalam keadaan~$q$
dan membaca simbol~$\sigma$, sedangkan fungsi transisinya tidak terdefinisi;
yakni, $M$ telah berhenti.
\end{proof}

\end{document}
